\subsection{A compact-perturbation lemma and the canonical heat branch}
\label{sec:short-seed-canonical}

We use the Gaussian angular-two Hilbert space $X_G$, its Euler operator
$L=T+A$ with $T=-2i\Omega$ and compact $A=2ik\Delta_2^{-1}$,
the simple eigenpair $(\lambda,q)$, and $P=q\ell$, $Q=I-P$ established
above. Write $\gamma=\Re\lambda>0$ and
\[
 D=\Delta+\tfrac12y\cdot\nabla+2,\qquad \kappa=\ell(Dq).
\]
On $X_G$, $D$ is selfadjoint and nonpositive; $q,q^\dagger$ belong to
every power domain of $D$. In particular $DP$ and $PD$ extend boundedly.
These statements concern the actual Gaussian adjoint, not an unweighted
formal adjoint.

\begin{lemma}[Compact Duhamel comparison]\label{lem:short-compact-duhamel}
Fix a finite time length $B$. Suppose uniformly bounded free evolutions
$F_n$ converge strongly to $F_0$, and their adjoints converge strongly,
uniformly in both times and any auxiliary parameters. Suppose the bounded
strongly measurable interactions $A_n(t)$ converge uniformly in operator
norm to $A_0(t)$,
and all the latter form a norm-precompact family of compact operators.
The full evolutions $U_n$ and their adjoints converge strongly, and
\[
 \sup_{0\le t-s\le B}
 \|(U_n-F_n)(t,s)-(U_0-F_0)(t,s)\|\longrightarrow0.
\]
\end{lemma}
\begin{proof}
The Dyson series have a uniform exponential majorant, so termwise
convergence proves the two strong assertions. For finite-rank
$K=\sum_{j=1}^m u_j\langle\,\cdot\,,v_j\rangle$, strong convergence
on $u_j$ and adjoint-strong convergence on $v_j$ imply norm convergence
of $F_n(t,r)KU_n(r,s)$. A finite norm-net and finite-rank approximation
of the compact family $A_0(r)$ make the conclusion uniform. Integrate
\[
 (U_n-F_n)(t,s)=\int_s^tF_n(t,r)A_n(r)U_n(r,s)\,dr.
\]
Uniform bounds absorb $A_n-A_0$. This proves norm convergence only for
the compact remainder, as required.
\end{proof}

\begin{lemma}[Small diffusion on the complementary space]
\label{lem:short-diffusion-complement}
There are $\beta_*,g,C>0$ such that, for every measurable
$0\le\beta(t)\le\beta_*$, the forward evolution on $QX_G$ generated by
\[
 Q\{L-\lambda+\beta(t)(D-\kappa)\}Q
\]
has norm at most $Ce^{-g(t-s)}$ for $t\ge s$.
\end{lemma}
\begin{proof}
Extend the generator to
$G_\beta=L-\lambda-P+\beta Q(D-\kappa)Q$ on $X_G$. Since
\[
 Q(D-\kappa)Q=D-\kappa+R_f,
 \quad R_f=-DP-PD+PDP+\kappa P,
\]
$R_f$ is bounded and finite rank. Split $G_\beta$ into the free
generator $T-\lambda-\beta\kappa+\beta D$ and interaction
$A-P+\beta R_f$. The evolution exists for measurable $\beta$ by the
Dyson series based on $\exp(D\int_s^t\beta)$; its short-interval
bounds are uniform. Dissipativity of $D$ and skew-adjointness of $T$
give
\[
 \|F_\beta(t,s)\|
 \le\exp\{-\gamma(t-s)+|\Re\kappa|\beta_*(t-s)\}.
\]
This inequality can first be proved for bounded nonpositive Yosida
approximants of $D$ and then passed to the strong limit.

On a fixed block, $\exp(D\int_s^t\beta)\to I$ strongly as
$\beta_*\downarrow0$, uniformly over all admissible coefficients.
The Dyson series give the corresponding direct and adjoint strong
limits for the free and full evolutions. For the adjoints, reverse
time on the block; the class of admissible coefficients is unchanged.
Lemma~\ref{lem:short-compact-duhamel} therefore gives
\[
 U_\beta(s+B,s)-F_\beta(s+B,s)
 \longrightarrow e^{B(L-\lambda-P)}-e^{B(T-\lambda)}
\]
in norm uniformly in $s$ and $\beta$. The inviscid complementary
semigroup estimate, with any rate below $\gamma$, makes
$L-\lambda-P$ exponentially stable. Consequently
\[
 \limsup_{\beta_*\downarrow0}\sup_{s,\beta}
 \|U_\beta(s+B,s)\|
 \le \|e^{B(L-\lambda-P)}\|+2e^{-\gamma B}<\tfrac14
\]
when $B$ is sufficiently large. Fixing $\beta_*$ then makes every
full block contract by $1/2$; iteration and the short-block bound
prove exponential decay. Finally $G_\beta$ commutes with $P$ and
equals $-1$ on its range. In the mild equation this follows by testing
against $q^\dagger\in\mathcal D(D)$ and uniqueness from the datum
$q$. Thus $QX_G$ is invariant and restriction proves the claim.
\end{proof}

The radial Gaussian evolves exactly by heat diffusion:
\[
 a=1+\nu s,\quad E(s,x)=a^{-2}E_0(x/\sqrt a),\quad
 V(s,x)=a^{-3/2}V_0(x/\sqrt a).
\]
It is also an unforced planar Navier--Stokes solution, since its
vorticity is radial and its velocity tangential. Substitution of
$w(s,x)=a^{-2}W(z,x/\sqrt a)$ gives the exact linearized equation
\begin{equation}\label{eq:short-heat-clock}
 W_z=(L+\theta^{-1}D)W,\qquad
 z=\frac{s}{1+\nu s}=\nu^{-1}-\theta,
 \qquad \theta=\frac1{\nu(1+\nu s)}.
\end{equation}
Here $dz/ds=a^{-2}$; the Biot--Savart velocity scales as $a^{-3/2}$.

\begin{proposition}[Canonical branch]
\label{prop:short-canonical}
There is a unique solution of
\[
 -V_{\rm can}'=(L+\theta^{-1}D)V_{\rm can},\qquad \theta>0,
\]
with
\begin{equation}\label{eq:short-canonical-asymptotic}
 V_{\rm can}(\theta)=e^{-\lambda\theta}\theta^{-\kappa}
     \{q+O_{X_G}(\theta^{-1})\}\quad(\theta\to\infty).
\end{equation}
It is bounded in $X_G$ as $\theta\downarrow0$ and smooth for positive
$\theta$. For $X_{G,j}=\mathcal D((I-D)^{j/2})$, fixed $j$, and
$\theta\ge b$ with a fixed sufficiently large $b$,
\[
 \|e^{\lambda\theta}\theta^\kappa V_{\rm can}(\theta)\|_{X_{G,j}}
       \le C_j(1+\theta)^{j/2}.
\]
\end{proposition}
\begin{proof}
Write $e^{\lambda\theta}\theta^\kappa V=Aq+r$, $r\in QX_G$.
The equations are
\[
 -A'=\theta^{-1}\ell(Dr),\qquad
 -r'=Q\{L-\lambda+\theta^{-1}(D-\kappa)\}Qr
                   +\theta^{-1}A\,QDq.
\]
For $b\ge\beta_*^{-1}$ let $S(\theta,s)$, $b\le\theta\le s$,
be the forward complementary evolution. It satisfies
$\|S(\theta,s)\|\le Ce^{-g(s-\theta)}$. The normalization at infinity
is imposed by
\[
 r_A(\theta)=\int_\theta^\infty
 S(\theta,s)QDq\,\frac{A(s)}s\,ds,\qquad
 (\mathcal TA)(\theta)=1+\int_\theta^\infty
                  \frac{\ell(Dr_A(s))}s\,ds.
\]
Because $\ell D$ is bounded,
$\|r_A(\theta)\|+|\mathcal TA(\theta)-1|
\le C\theta^{-1}\|A\|_\infty$. The same estimate for differences
makes $\mathcal T$ a contraction on $C_b([b,\infty))$ after enlarging
$b$. Its fixed point gives \eqref{eq:short-canonical-asymptotic}.
Splitting both integrals at any finite $s_0$ verifies the coupled
mild initial-value equations, whose uniqueness follows from the
bounded coupling maps and Gronwall. Conversely every normalized
solution satisfies the integrals: the term
$S(\theta,s_0)r(s_0)$ tends to zero as $s_0\to\infty$.
This also proves uniqueness at infinity.

Set $\tau=\log(b/\theta)$. Then
\[
 \partial_\tau V=DV+be^{-\tau}LV.
\]
The bounded-perturbation formula and $D\le0$ give continuation for
all $\tau\ge0$ with
$\sup_{0<\theta\le b}\|V(\theta)\|
\le e^{b\|L\|}\|V(b)\|$. This is forward parabolic continuation;
no backward heat inverse occurs.

For completeness, conjugation by $G^{1/2}$ turns $D$ into a shifted
harmonic oscillator. Its Hermite decomposition shows that the squared
$X_{G,j}$ norm is equivalent to
$\sum_{|\alpha|\le j}\|\partial^\alpha f\|_{L^2(G^{-1})}^2$
(Cartesian derivatives may leave angular sector two). Multiplication
by $\Omega$ is bounded on these spaces. The radial Green formula for
$\Delta_2^{-1}$ first gives $\|\Delta_2^{-1}f\|_2\le C\|f\|_{X_G}$;
elliptic regularity and boundedness of
$G^{-1/2}\partial^\alpha k$ then show $L:X_{G,j}\to X_{G,j}$ is
bounded. Spectral calculus gives a one-derivative heat gain
$\|e^{aD}\|_{X_{G,j}\to X_{G,j+1}}\le C_ja^{-1/2}$ for $0<a\le1$.
On the unit turnover interval from $\theta+1$ down to $\theta$,
diffusion is comparable to $(1+\theta)^{-1}$. Duhamel therefore gains
one derivative at elapsed time $h$ at cost
$C_j((1+\theta)/h)^{1/2}$; its integral kernel is integrable.
Divide the interval into $j$ equal pieces to obtain the stated
$C_j(1+\theta)^{j/2}$ bound from the bounded $X_G$ norm at its inlet.
The $\tau$ equation propagates every fixed graph norm to zero because
$\int_0^\infty be^{-\tau}\,d\tau=b$.
\end{proof}

\begin{corollary}[Nontrivial full-age output]\label{cor:short-mellin}
Let $\mathfrak M(f)=\int_0^\infty r^3f_2(r)\,dr$ and
$M_q=\mathfrak M(q)\ne0$. Then
\[
 \mathcal J(\mu)=\int_0^\infty
       \theta^{\mu+1/2}\mathfrak M(V_{\rm can}(\theta))\,d\theta
\]
is holomorphic and not identically zero on $\Re\mu>-3/2$.
\end{corollary}
\begin{proof}
Weighted Cauchy--Schwarz bounds $\mathfrak M$ on $X_G$. Thus its
canonical value is $O(1)$ at zero and
$e^{-\lambda\theta}\theta^{-\kappa}(M_q+O(\theta^{-1}))$ at infinity.
These bounds prove absolute convergence and permit differentiation
under the integral, including every logarithmic factor on compact
subsets of the half-plane. If $\mathcal J$ vanished identically, then
for fixed real $b>-3/2$ the $L^1(\mathbb R)$ function
$e^{(b+3/2)x}\mathfrak M(V_{\rm can}(e^x))$ would have identically zero
Fourier transform. Fourier uniqueness would contradict $M_q\ne0$.
\end{proof}
